En esta evaluación se dispone de AUC para 12 de los doce hallazgos. El AUC macro de doce etiquetas de las variantes varía entre 0.513 y 0.620; F1 macro varía entre 0.361 y 0.506. Estos intervalos describen la matriz ejecutada y no son intervalos de confianza. No se elige una variante para despliegue por estos resultados de evaluación. Las curvas muestran el criterio aplicado para elegir cada checkpoint; no se comparan los valores absolutos de BCE y focal como si fueran la misma escala.  Comparaciones pareadas por arquitectura y tratamiento del ruido: CNN/bce/pond.: diferencia focal menos BCE de AUC -0.005 y de F1 macro +0.135; CNN/bce/filtro: diferencia focal menos BCE de AUC +0.007 y de F1 macro -0.028; ResNet/bce/pond.: diferencia focal menos BCE de AUC +0.001 y de F1 macro +0.144; ResNet/bce/filtro: diferencia focal menos BCE de AUC -0.014 y de F1 macro +0.035. Con una semilla y una cohorte experta pequeña, estas diferencias no establecen superioridad general. La mitigación del desbalance debe evaluarse por hallazgo y soporte, incluyendo posibles pérdidas de precisión al elevar recall. El umbral 0,5 permanece fijo aun cuando una variante produce probabilidades mal calibradas. La distribución tras el filtro explica qué etiquetas conservan positivos y negativos explícitos; no se interpreta el filtrado como limpieza perfecta.